\begin{mistake}{2026-04-12}{16}

\begin{problemcontent}
求
\[
\int_{-1}^{0} \frac{\ln(1+x)}{\sqrt[3]{1+x}}\,dx.
\]

\end{problemcontent}

\begin{solutioncontent}
令 \(t=1+x\)，则 \(dx=dt\)。当 \(x=-1\) 时 \(t=0\)；当 \(x=0\) 时 \(t=1\)。原积分化为
\[
\int_0^1 t^{-1/3}\ln t\,dt.
\]

这是瑕积分，应写成
\[
\int_0^1 t^{-1/3}\ln t\,dt
=\lim_{\varepsilon\to 0^+}\int_{\varepsilon}^1 t^{-1/3}\ln t\,dt.
\]

对积分作分部积分。取
\[
u=\ln t,
\qquad dv=t^{-1/3}\,dt,
\]
则
\[
du=\frac{1}{t}\,dt,
\qquad v=\frac{3}{2}t^{2/3}.
\]

于是
\[
\begin{aligned}
\int t^{-1/3}\ln t\,dt
&=\frac{3}{2}t^{2/3}\ln t-\frac{3}{2}\int t^{-1/3}\,dt\\
&=\frac{3}{2}t^{2/3}\ln t-\frac{3}{2}\cdot\frac{3}{2}t^{2/3}+C\\
&=\frac{3}{2}t^{2/3}\ln t-\frac{9}{4}t^{2/3}+C.
\end{aligned}
\]

故原积分为
\[
\lim_{\varepsilon\to 0^+}\left[\frac{3}{2}t^{2/3}\ln t-\frac{9}{4}t^{2/3}\right]_{\varepsilon}^{1}.
\]

其中
\[
\lim_{t\to 0^+} t^{2/3}\ln t
=\lim_{t\to 0^+}\frac{\ln t}{t^{-2/3}}
=\lim_{t\to 0^+}\frac{1/t}{-(2/3)t^{-5/3}}
=\lim_{t\to 0^+}\left(-\frac{3}{2}t^{2/3}\right)=0.
\]

又有 \(\lim_{t\to 0^+} t^{2/3}=0\)，故
\[
\int_{-1}^{0} \frac{\ln(1+x)}{\sqrt[3]{1+x}}\,dx
=\left(0-\frac{9}{4}\right)-0
=-\frac{9}{4}.
\]

\end{solutioncontent}

\mistakeknowledge{
\begin{itemize}
  \item \textbf{核心公式/定理}：换元后先识别瑕积分，再用分部积分处理 \(\int_0^1 t^{-1/3}\ln t\,dt\)。其中 \(\ln t\) 优先取为 \(u\)。
  \item \textbf{易错点}：不能把下限 \(0\) 直接代入原函数，必须先写成极限；否则对 \(t^{2/3}\ln t\) 的处理容易失去严格性。
  \item \textbf{关联知识}：\(\lim_{t\to 0^+} t^\alpha\ln t=0\,(\alpha>0)\) 是常见极限结论，可由洛必达法则或指数替换证明。
\end{itemize}}

\end{mistake}
